% GENERATED FILE. Edit canonical lessons, then run npm run generate:books.
% canonical-source-hash: fnv1a64:46b99a6b7079b73c
% canonical-lessons: FR-C132-bruyant, FR-C132-dernier, FR-C132-prochain, FR-C132-celebre, FR-C132-drole

\chapter{Noisy, Last, Next, Famous, Funny}
\label{ch:fr-noisy-last-next-famous-funny}
\hlchaptermodality{132}

\begin{chapteropening}
\textbf{By the end of this chapter:} \emph{I can name noisy, last, next, famous, funny.}

Complete the last lesson of chapter 132: I can name noisy, last, next, famous, funny.
\end{chapteropening}

\section[bruyant]{bruyant — noisy}

\label{lesson:FR-C132-bruyant}



\begin{quote}
Before the new one: say the French for closed, then the French for calm.
\end{quote}

\subsection*{You'll want to know: bruyant}
\textbf{bruyant} — \textquotedblleft{}noisy\textquotedblright{}.

\begin{cousinweb}
\textbf{bruyant} is the present participle of the Old French verb \emph{bruire}, \textquotedblleft{}to roar, to rustle\textquotedblright{}, usually traced to Latin \emph{rugire}, \textquotedblleft{}to roar\textquotedblright{} (Lewis \& Short), with a changed start.
\end{cousinweb}

\begin{grammarlens}[title={What you've built}]
One more word for this chapter.
\end{grammarlens}

\subsection*{Your turn}
\begin{itemize}
\raggedright
  \item \emph{Say it:} \emph{bruyant}
  \item \emph{Say it:} \emph{bruyant}, once more
  \item \emph{Say it:} say \emph{bruyant}
  \item \emph{Recall:} say the French for closed, then the French for calm, then say \emph{bruyant} again
\end{itemize}

\begin{tcolorbox}[breakable,colback=teal!4,colframe=teal!35!black,title={Before you move on}]
What does bruyant mean? (\textquotedblleft{}Noisy\textquotedblright{}.) Say it once more.
\end{tcolorbox}
\section[dernier]{dernier — last}

\label{lesson:FR-C132-dernier}



\begin{quote}
Before the new one: say the French for calm, then the French for noisy.
\end{quote}

\subsection*{You'll want to know: dernier}
\textbf{dernier} — \textquotedblleft{}last\textquotedblright{}.

\begin{cousinweb}
\textbf{dernier} — built in Old French on \emph{derrain}, \textquotedblleft{}last\textquotedblright{}, from Vulgar Latin \emph{*deretranus}, built on Latin \emph{retro}, \textquotedblleft{}behind\textquotedblright{} (Lewis \& Short).
\end{cousinweb}

\begin{grammarlens}[title={What you've built}]
One more word for this chapter.
\end{grammarlens}

\subsection*{Your turn}
\begin{itemize}
\raggedright
  \item \emph{Say it:} \emph{dernier}
  \item \emph{Say it:} \emph{dernier}, once more
  \item \emph{Say it:} say \emph{dernier}
  \item \emph{Recall:} say the French for calm, then the French for noisy, then say \emph{dernier} again
\end{itemize}

\begin{tcolorbox}[breakable,colback=teal!4,colframe=teal!35!black,title={Before you move on}]
What does dernier mean? (\textquotedblleft{}Last\textquotedblright{}.) Say it once more.
\end{tcolorbox}
\section[prochain]{prochain — next}

\label{lesson:FR-C132-prochain}



\begin{quote}
Before the new one: say the French for noisy, then the French for last.
\end{quote}

\subsection*{You'll want to know: prochain}
\textbf{prochain} — \textquotedblleft{}next\textquotedblright{}.

\begin{cousinweb}
\textbf{prochain} $\leftarrow$ Vulgar Latin \emph{*propeanus}, from Latin \emph{prope}, \textquotedblleft{}near\textquotedblright{} — inherited. (EtymDB 2.1)
\end{cousinweb}

\begin{grammarlens}[title={What you've built}]
One more word for this chapter.
\end{grammarlens}

\subsection*{Your turn}
\begin{itemize}
\raggedright
  \item \emph{Say it:} \emph{prochain}
  \item \emph{Say it:} \emph{prochain}, once more
  \item \emph{Say it:} say \emph{prochain}
  \item \emph{Recall:} say the French for noisy, then the French for last, then say \emph{prochain} again
\end{itemize}

\begin{tcolorbox}[breakable,colback=teal!4,colframe=teal!35!black,title={Before you move on}]
What does prochain mean? (\textquotedblleft{}Next\textquotedblright{}.) Say it once more.
\end{tcolorbox}
\section[célèbre]{célèbre — famous}

\label{lesson:FR-C132-celebre}



\begin{quote}
Before the new one: say the French for last, then the French for next.
\end{quote}

\subsection*{You'll want to know: célèbre}
\textbf{célèbre} — \textquotedblleft{}famous\textquotedblright{}.

\begin{cousinweb}
\textbf{célèbre} $\leftarrow$ Latin \emph{celeber}, \textquotedblleft{}crowded, much spoken of\textquotedblright{} — a learned borrowing. (EtymDB 2.1)
\end{cousinweb}

\begin{grammarlens}[title={What you've built}]
One more word for this chapter.
\end{grammarlens}

\subsection*{Your turn}
\begin{itemize}
\raggedright
  \item \emph{Say it:} \emph{célèbre}
  \item \emph{Say it:} \emph{célèbre}, once more
  \item \emph{Say it:} say \emph{célèbre}
  \item \emph{Recall:} say the French for last, then the French for next, then say \emph{célèbre} again
\end{itemize}

\begin{tcolorbox}[breakable,colback=teal!4,colframe=teal!35!black,title={Before you move on}]
What does célèbre mean? (\textquotedblleft{}Famous\textquotedblright{}.) Say it once more.
\end{tcolorbox}
\section[drôle]{drôle — funny}

\label{lesson:FR-C132-drole}



\begin{quote}
Before the new one: say the French for next, then the French for famous.
\end{quote}

\subsection*{You'll want to know: drôle}
\textbf{drôle} — \textquotedblleft{}funny\textquotedblright{}.

\begin{cousinweb}
\textbf{drôle} $\leftarrow$ Middle Dutch \emph{drol}, \textquotedblleft{}a little fellow, a goblin\textquotedblright{} — borrowed; it first meant \textquotedblleft{}a rascal\textquotedblright{}, then \textquotedblleft{}a comic\textquotedblright{}.
\end{cousinweb}

\begin{grammarlens}[title={What you've built}]
One more word for this chapter.
\end{grammarlens}

\subsection*{Your turn}
\begin{itemize}
\raggedright
  \item \emph{Say it:} \emph{drôle}
  \item \emph{Say it:} \emph{drôle}, once more
  \item \emph{Say it:} say \emph{drôle}
  \item \emph{Recall:} say the French for next, then the French for famous, then say \emph{drôle} again
\end{itemize}

\begin{tcolorbox}[breakable,colback=teal!4,colframe=teal!35!black,title={Before you move on}]
What does drôle mean? (\textquotedblleft{}Funny\textquotedblright{}.) Say it once more.
\end{tcolorbox}
